% 03-rnn-don-gian.tex - The plain RNN forward pass
%
% Scope: forward equations with pre-activation; unrolled view; sharing
% weights across time; the loss function.

\section{RNN: lan truyền xuôi}
\label{sec:ch01-rnn-xuoi}

Với ký hiệu đã có, tôi viết lại lan truyền xuôi của RNN một cách tường minh
cho toàn bộ \tn{chuỗi}{sequence}. Đây không phải là chỗ để đọc lướt; mỗi ký
hiệu trong dẫn
xuất BPTT ở mục~4 đều quay về một dòng ở đây.

\subsection{Một bước thời gian}

Tại bước $t$, ba thứ xảy ra theo thứ tự:

\begin{enumerate}
  \item Tính pre-activation:
    $\vect{a}_t = \mat{W}_{xh} \vect{x}_t
                 + \mat{W}_{hh} \vect{h}_{t-1}
                 + \vect{b}_h$.
  \item Qua \tn{hàm kích hoạt}{activation function}:
    $\vect{h}_t = \tanh(\vect{a}_t)$,
    với $\tanh$ áp dụng lên từng phần tử.
  \item Đọc ra output:
    $\vect{y}_t = \mat{W}_{hy} \vect{h}_t + \vect{b}_y$.
\end{enumerate}

Cùng một bộ trọng số $(\mat{W}_{xh}, \mat{W}_{hh}, \mat{W}_{hy},
\vect{b}_h, \vect{b}_y)$ được dùng ở mọi bước. Đây là \textbf{chia sẻ trọng
số} (weight sharing): RNN không học một bộ tham số riêng cho $t=1$, một bộ
khác cho $t=2$, mà học một bộ duy nhất áp dụng cho chuỗi dài bất kỳ. Nhờ vậy,
một RNN huấn luyện trên câu dài 10 từ vẫn chạy được trên câu dài 50 từ mà
không cần thay đổi gì.

\subsection{Toàn bộ chuỗi}

Với chuỗi dài $T$, lan truyền xuôi sinh ra ba dãy:

\[
  \vect{h}_1, \vect{h}_2, \dots, \vect{h}_T; \qquad
  \vect{y}_1, \vect{y}_2, \dots, \vect{y}_T; \qquad
  \vect{a}_1, \vect{a}_2, \dots, \vect{a}_T.
\]

\tn{Trạng thái ẩn}{hidden state} ban đầu $\vect{h}_0 = \vect{0}$ (vector
không). Đây là lựa chọn
đơn giản nhất và là mặc định trong mọi cài đặt; sau này LSTM và GRU cũng khởi
tạo cell state và hidden state bằng vector không.

\subsection{Hàm mất mát}

Tôi dùng \textbf{tổng bình phương sai số} trên toàn bộ chuỗi. Với
$\widehat{\vect{y}}_t$ là target tại bước $t$:

\begin{equation}
  L = \frac{1}{2} \sum_{t=1}^{T} \bigl\|
      \vect{y}_t - \widehat{\vect{y}}_t
      \bigr\|^2.
  \label{eq:rnn-loss}
\end{equation}

Hệ số $1/2$ không thay đổi vị trí cực tiểu nhưng làm gọn gradient: đạo hàm
của $z^2/2$ là $z$, không phải $2z$. Trong thực tế, bài toán phân loại chuỗi
dùng cross-entropy, và bài toán sinh chuỗi dùng cross-entropy trên từng
token. Nhưng những \tn{hàm mất mát}{loss function} đó thêm một tầng softmax và
một tầng Jacobian
không giúp ích gì cho việc hiểu BPTT. Tổng bình phương sai số giữ gradient
ở dạng đơn giản nhất: một phép trừ.

Chương này chỉ dùng bài toán hồi quy; dự đoán một vector liên tục từ một
chuỗi đầu vào. Đó là bài toán nhân tạo, nhưng mọi thứ ta học từ nó áp dụng
nguyên vẹn cho bài toán thật. BPTT không quan tâm bạn đặt gì sau
$\vect{y}_t$; nó chỉ cần biết $\partial L / \partial \vect{y}_t$, và con số
đó đến từ đâu không ảnh hưởng đến phần còn lại của dẫn xuất.
